\clearpage\SourceLeaf{314}

\begin{center}
[POTHENOTS AUFGABE UND DAS VIERECK.]
\end{center}

\begin{center}
[I.]

[CONSTRUCTION DER POTHENOTSCHEN AUFGABE.]

[1.]
\end{center}

Es werden drei Punkte $a$, $b$, $c$ durch die eben so bezeichneten complexen
Zahlen dargestellt, es seien ferner $\alpha$, $\beta$, $\gamma$ drei complexe
Zahlen, so dass
\[
\alpha+\beta+\gamma=0.
\]
Man mache
\[
\alpha A+\beta b+\gamma c=0
\]
\[
\alpha a+\beta B+\gamma c=0
\]
\[
\alpha a+\beta b+\gamma C=0,
\]
und $A$, $B$, $C$ seien die Punkte, denen die complexen Zahlen
$A$, $B$, $C$ entsprechen. Dann schneiden sich $aA$,
$bB$, $cC$ in Einem Punkte $P$ und dieser ist die Auflösung
von \textsc{Pothenots} Aufgabe.

Sind $\mathfrak{A}$, $\mathfrak{B}$, $\mathfrak{C}$ die in $P$ beobachteten
(nicht orientirten) Azimuthe, so verhalten sich
\[
\alpha e^{i\mathfrak{A}}, \qquad \beta e^{i\mathfrak{B}}, \qquad \gamma e^{i\mathfrak{C}}.
\]
wie reelle Zahlen.

\Needspace{8\baselineskip}
\begin{center}
BEMERKUNG.
\end{center}

Diese Notiz hat \textsc{Gauss} auf der innern Seite des hintern Deckels
seines Exemplars des Werkes: \textsc{Carnot}, \emph{Geometrie der
Stellung}. Übersetzt von \textsc{Schumacher}. Bd.\ II.\ Altona 1810
eingetragen.

\par\hfill\textsc{Stäckel}.\par
